% COMSYS Seminar/Preseminar Paper Template
% ----------------------------------------
% Based on the original IEEE template package, includes annotations by COMSYS

\documentclass[conference]{IEEEtran}

%\IEEEoverridecommandlockouts
\usepackage{amsmath,amssymb,amsfonts}
\usepackage{algorithmic}
\usepackage[hidelinks]{hyperref}
\hypersetup{hypertexnames=false}
\usepackage[capitalize,nameinlink]{cleveref}
\usepackage[sort]{cite}
\usepackage[inline]{enumitem}
\usepackage{graphicx}
\usepackage{textcomp}
\usepackage{array}
\usepackage{tabularx}
\usepackage{booktabs}
\usepackage{float}
\usepackage{xcolor}
\usepackage{listings}
\def\BibTeX{{\rm B\kern-.05em{\sc i\kern-.025em b}\kern-.08em
    T\kern-.1667em\lower.7ex\hbox{E}\kern-.125emX}}

\definecolor{listingbg}{RGB}{248,248,248}
\definecolor{listingrule}{RGB}{185,185,185}
\definecolor{listingnumber}{RGB}{100,100,100}
\lstdefinelanguage{MIR}{
  morekeywords={param,goto,return,ifTrue,ifFalse},
  sensitive=true,
  morecomment=[l]{//},
  morestring=[b]"
}
\lstdefinelanguage{Pseudo}{
  morekeywords={algorithm,for,each,every,if,else,return,repeat,until,while,where,with,from,to,in,not,and,or,continue,then},
  sensitive=false,
  morecomment=[l]{//},
  morestring=[b]"
}
\lstdefinestyle{paperlisting}{
  basicstyle=\ttfamily\scriptsize,
  frame=tb,
  framerule=0.4pt,
  columns=fullflexible,
  keepspaces=true,
  showstringspaces=false,
  breaklines=true,
  breakatwhitespace=false,
  tabsize=2,
  xleftmargin=0em,
  xrightmargin=0em,
  aboveskip=0.6\baselineskip,
  belowskip=0.6\baselineskip,
  captionpos=b
}
\lstdefinestyle{numberedlisting}{
  style=paperlisting,
  numbers=left,
  numberstyle=\tiny\color{listingnumber},
  numbersep=5pt,
  xleftmargin=1.6em,
  framexleftmargin=1.6em
}
\lstdefinestyle{cfull}{style=numberedlisting,language=C}
\lstdefinestyle{mirfull}{style=numberedlisting,language=MIR}
\lstdefinestyle{csmall}{style=paperlisting,language=C}
\lstdefinestyle{mir}{style=paperlisting,language=MIR}
\lstdefinestyle{pseudo}{style=paperlisting,language=Pseudo}
\crefname{lstlisting}{listing}{listings}
\Crefname{lstlisting}{Listing}{Listings}

\newcolumntype{Y}{>{\raggedright\arraybackslash}X}
\newcommand{\tighttable}{\scriptsize\renewcommand{\arraystretch}{1.16}}

% Treat the "if-then-else" selector as a single function-name token rather than
% three italic letters, so it does not read like a product of variables.
\newcommand{\ite}{\mathit{ite}}

\begin{document}

\title{Explaining State-Merging Techniques Through MIR Dumps}

\author{\IEEEauthorblockN{Yaroslav Riabtsev}
\IEEEauthorblockA{\textit{ Communication and Distributed Systems} \\
\textit{RWTH Aachen University}\\
Aachen, Germany \\
yaroslav.riabtsev@rwth-aachen.de}
}

\maketitle

\begin{abstract}
This article explains state merging through inspectable MIR examples and dumps.
Instead of evaluating a full symbolic executor by coverage, bugs found,
generated tests, or solver runtime, it follows the intermediate structures that
make merging decisions understandable: path splits and joins, conditional
symbolic values, dependency sets, future-relevant variables, and compact
summaries of repeated or regional control flow.

The paper starts from dynamic symbolic execution and standard state merging,
where merging reduces the number of states but moves complexity into
disjunctions and \(\ite\)-expressions. It then uses small MIR examples to
explain Query Count Estimation, hot variables, state similarity, Dynamic State
Merging, Veritesting-style region summarization, and indexed loop summaries for
regular path families.

The central theme is the same throughout: a merge is useful when the
formula complexity it introduces remains manageable. The MIR dumps make this
trade-off visible by showing which path-dependent values reach future
conditions, which variables become important for later queries, and which CFG
or loop structures admit a compact summary.
\end{abstract}


\begin{IEEEkeywords}
symbolic execution, dynamic symbolic execution, state merging, query count estimation, dynamic state merging, Veritesting, static symbolic execution, MIR, control-flow graph, path explosion
\end{IEEEkeywords}

\section{Introduction}
\label{sec:introduction}

KLEE is one of the most widely known practical systems for dynamic symbolic
execution. It takes programs represented in LLVM IR, executes them over
symbolic inputs, and maintains symbolic states corresponding to different
possible executions of the program. This execution model is useful in automated
test generation, coverage-oriented exploration, bug finding, and reachability
analysis. For this reason, KLEE provides a convenient technical reference point
for discussing more specialized symbolic-analysis techniques \cite{klee}.

Several techniques considered in this article can be understood as extensions,
modifications, or close relatives of a KLEE-like execution model. They preserve
the same general view of a program as executable intermediate code with explicit
control flow, but change how symbolic states, path predicates, program regions,
or repeated path structures are represented and manipulated. The article
therefore looks inside these techniques rather than treating them as black boxes,
focusing on the local mechanics: which data an algorithm needs, which
structures appear in the CFG, which expressions are constructed, which
heuristics guide decisions, and which intermediate artifacts determine these
steps.

Throughout the article, a symbolic state is written as

\[
s = (\ell, \pi, \mu).
\]

Here, \(\ell\) is the current program location: an instruction, a basic block,
or a point in the CFG. The component \(\pi\) is the path condition, i.e. the
constraint accumulated along the current path. The component \(\mu\) is the
symbolic store: a mapping from program variables to symbolic expressions.
Accordingly, \(\mu(x)\) denotes the current symbolic value of variable \(x\).
Different papers use different notation for this conceptual triple, such as
location/path condition/store, instruction counter/path predicate/symbolic
memory, or variants like \((\ell, \Pi, \Delta)\). This article uses
\((\ell, \pi, \mu)\) consistently.

KLEE operates on LLVM IR, but the experiments in this article use a smaller
MIR-like representation. The representation contains only the constructs needed
to demonstrate the techniques under discussion: parameters, assignments,
arithmetic operations, conditional jumps, unconditional jumps, and returns. This
level is sufficient to expose basic blocks, CFGs, join points, branch
conditions, simple loops, variable dependencies, and candidate summaries.

The examples and tables were produced using an existing MIR analysis project
reused as a lightweight playground \cite{compilers-course}. The MIR experiments
inspect static control-flow and dependency information; they do not perform
path-feasibility solving or test generation. The article itself is
self-contained: the dumps provide reproducibility artifacts, but the relevant
listings, equations, and result tables are included in the text.

\section{Motivation and Baseline}
\label{sec:motivation}

Symbolic execution occupies a useful position between concrete execution and
static reasoning about program behavior. In a concrete run, a program receives
specific input values and follows one specific path. In symbolic execution,
some inputs are replaced with symbolic values, and program operations construct
expressions over those symbols. As a result, a single analysis state can
represent not just one execution, but a family of executions satisfying a
particular condition.

The practical motivation for such systems is usually described through the
outputs of the analysis. If the condition for following a path is satisfiable,
a solver can construct a concrete input that makes the original program follow
the same path. This makes symbolic execution useful for automatic test
generation, line and branch coverage improvement, bug finding, reachability
checking, and reproducible bug reports. This connection to concrete inputs is
one of the reasons symbolic execution remains attractive: an analysis result
can often be turned into an actual program run.

Another important layer is the SMT-solving layer. In practical symbolic
execution systems, path conditions and symbolic expressions over inputs are
translated into logical formulas. A solver is then used to check path
feasibility, discard impossible branches, and construct models for concrete
inputs. This layer accounts for a significant part of the practical cost of
the analysis: even when the number of explored paths is moderate, the formulas
themselves may become difficult for the solver.

This article focuses on the algorithmic foundations that support both of these
aspects. Test generation, coverage, verification, and bug reporting provide the
external motivation rather than experimental metrics. Instead, the article
studies the intermediate structures on which such systems rely: symbolic states,
path conditions, symbolic stores, control-flow graphs, join points, conditional
expressions, variable dependencies, and local indicators of future complexity.

This level is especially useful for discussing state-merging techniques. These
techniques are not reducible to a single formula or a single solver call. They
use information from different parts of the analysis: where join points are
located, which values differ across paths, whether those values may appear in
future conditions, whether several paths can be represented by a regional
summary, and whether a repeated pattern is present in a loop. Understanding
the algorithms therefore requires looking at the intermediate artifacts on
which their decisions are based.

\subsection{Dynamic Symbolic Execution}
\label{sec:dse}

KLEE is a path-based symbolic execution engine that dynamically interprets LLVM IR over symbolic states. A program is executed step
by step, but input values may be symbolic rather than concrete. Operations over
symbolic values construct expressions, and conditional branches add constraints
to the path condition of the current state.

Consider the example in \Cref{lst:c-count-matches}:

\begin{lstlisting}[style=cfull,caption={C loop with repeated symbolic branch.},label={lst:c-count-matches}]
int count_matches(int i, int n, int target) {
    int matches = 0;

    while (i < n) {
        if (target == i) {
            matches = matches + 1;
        } else {
            matches = matches + 0;
        }

        i = i + 1;
    }

    return matches;
}
\end{lstlisting}

Assume that \(i\), \(n\), and \(\mathit{target}\) are symbolic inputs. At line 4, the
executor has to determine whether the loop continues. At line 5, it has to
choose one of the two updates of \texttt{matches}.

When a branch condition depends on symbolic data, KLEE performs a fork: the
current state is cloned, and the analysis continues with two states. One state
receives the constraint for the true branch, and the other receives the
constraint for the false branch. For example, entering the loop and taking the
true branch of the inner condition gives the path condition
\(i < n \land \mathit{target} = i\). Taking the false branch instead produces another
state with the condition \(i < n \land \lnot(\mathit{target} = i)\).

After the update \texttt{i = i + 1}, control returns to the loop condition, and
the same branch structure repeats for the next value of \(i\). If the loop is
unrolled for \(k\) iterations and both continuations of the inner branch are
kept as syntactic variants, then the inner branch alone defines up to \(2^k\)
branch histories.

For this concrete scalar example, an SMT-backed executor may prune many
combinations, because \(\mathit{target} = i\) can be true for at most one loop
iteration. The MIR-level experiments keep the syntactic path family induced by
the CFG; this gives the structural path family before feasibility pruning.

KLEE does not add such branches blindly. For each side of a conditional branch,
it checks the satisfiability of the corresponding path condition. If
\(\pi \land c\) is unsatisfiable, the true branch is discarded. If
\(\pi \land \lnot c\) is unsatisfiable, the false branch is discarded. The SMT
solver therefore partially limits the growth of states: a syntactically
available branch may be infeasible and will not be explored further.

The experiments use the MIR-like version of the same example, shown in
\Cref{lst:mir-count-matches}.

\begin{lstlisting}[style=mirfull,caption={MIR form of the repeated-branch loop.},label={lst:mir-count-matches}]
param i
param n
param target
matches <-- #0
Lloop: ifFalse i < n goto Lend
ifFalse target = i goto Lskip
matches <-- +, matches, #1
goto Lnext
Lskip: matches <-- +, matches, #0
Lnext: i <-- +, i, #1
goto Lloop
Lend: return matches
\end{lstlisting}

The MIR form makes basic blocks and transitions explicit: the loop condition
becomes a conditional jump from \texttt{Lloop}, the inner condition becomes a
second conditional jump, and the update of \texttt{i} returns control to the
loop header. This structure is enough to inspect CFG-level path growth: if the
CFG contains two syntactic continuations, both are visible as potential
variants in the structural experiment.

As a result, the experiment exposes path growth in its pure structural form.
Each binary symbolic branch doubles the number of potential histories, and a
repeated branch inside a loop produces an exponential family of paths. This is
path explosion.

\subsection{Standard State Merging}
\label{sec:state-merging}

Path explosion arises because dynamic symbolic execution keeps different
branch histories as different states. If paths diverge permanently, this is
natural: they describe genuinely different continuations of the execution.
However, in many programs, branches later reconverge at the same program
location. In that case, it becomes possible to continue the analysis not from
several independent states, but from one merged state.

The idea of standard state merging is to combine several states located at the
same program point \cite{efficient-merging,multise}. For simplicity, the
following equations assume that the two states are
merge-compatible: they are located at the same program point and their stores
are defined over the same variables relevant after the merge.

Let two such states be

\[
s_1 = (\ell, \pi_1, \mu_1),
\]

\[
s_2 = (\ell, \pi_2, \mu_2).
\]

Both states are at the same location \(\ell\), but they reached it through
different paths. The merged state must preserve information about both paths.
Therefore, its path condition becomes a disjunction:

\begin{equation}
\label{eq:merged-path-condition}
\pi = \pi_1 \lor \pi_2.
\end{equation}

Variable values must also depend on the path from which the execution arrived.
For this purpose, state merging uses conditional expressions of the form
\(\ite\), meaning if-then-else. For a variable \(x\), the merged store can be
written as

\begin{equation}
\label{eq:merged-store}
\mu(x) = \ite(\pi_1, \mu_1(x), \mu_2(x)).
\end{equation}

That is, if the path condition of the first state holds, the value is taken
from \(\mu_1(x)\); otherwise, it is taken from \(\mu_2(x)\). In general,
merging reduces the number of states, but moves path differences into formulas.

In the example from the previous subsection, this can be seen on the variable
\texttt{matches}. In \Cref{lst:c-count-matches}, lines 5--9 select one of two
updates: line 6 increments \texttt{matches}, while line 8 leaves it unchanged.

After these branches, control reaches the same point before the update of
\texttt{i}. Instead of keeping two separate states, the value of
\texttt{matches} can be represented by one conditional expression:

\[
\mathit{matches}' =
\ite(\mathit{target} = i,\ \mathit{matches} + 1,\ \mathit{matches} + 0).
\]

This representation follows the store shape in \Cref{eq:merged-store} and
remains precise: both branches are still represented in the formula. At the
same time, the number of active states decreases, because the later analysis
can continue from one merged state.

However, standard state merging does not remove complexity; it changes its
form. Before merging, the complexity is visible as the number of states and
paths. After merging, it appears inside the path condition and symbolic store:
disjunctions and \(\ite\)-expressions are introduced. If such conditional values
are later used in new branch conditions, solver queries may become harder.

Thus, state merging is a trade-off. It can reduce path explosion at the level of
the number of states, but it may increase formula complexity. The following
parts of the article examine this trade-off: when merging looks useful, when it
becomes risky, and which intermediate data helps make that decision.

\paragraph{ITE-cost shape: conditional values in later branches}
To make the cost of state merging visible, it is useful to look both at the
fact that states are merged and at where the introduced \(\ite\)-expression
flows afterwards. If a conditional value is not used in later branch
conditions, the merge may look relatively safe. If the value reaches a future
branch condition, then the complexity introduced by merging is pushed into a
later query.

Consider the MIR fragment in \Cref{lst:mir-ite-preview}:

\begin{lstlisting}[style=mirfull,caption={MIR branch with a join and a later hot branch.},label={lst:mir-ite-preview}]
param x
param y
ifTrue x < #0 goto Lneg
a <-- +, x, #5
goto Ljoin
Lneg: neg <-- -, #0, x
a <-- +, neg, #1
Ljoin: t <-- +, a, y
ifTrue t > #0 goto Lhot
r <-- +, y, #1
goto Lend
Lhot: r <-- +, t, #1
Lend: return r
\end{lstlisting}

The first branch splits execution on the condition \(x < 0\).

At \texttt{Ljoin}, the variable \texttt{a} has different expressions depending
on the selected branch:

\[
a =
\begin{cases}
1 - x, & x < 0,\\
x + 5, & \lnot(x < 0).
\end{cases}
\]

A state-merge-style representation of this value can be written as
\(a = \ite(x < 0,\ 1 - x,\ x + 5)\). Immediately after the join block,
the program computes \(t = a + y\), and then checks the condition \(t > 0\).
The incoming path conditions are mutually exclusive, and the conditional value
is interpreted under their disjunction.

Substituting the shallow \(\ite\)-representation of \texttt{a} into this branch
condition gives

\[
\ite(x < 0,\ 1 - x,\ x + 5) + y > 0.
\]

This is the important shape for the later analysis: the \(\ite\)-expression
does not simply sit in the symbolic store; it reaches a future branch condition.
That makes the merge candidate interesting and potentially risky.

For the merge site \texttt{Ljoin} and the variable \texttt{a}, the experiment
reconstructs the sequence
\[
  a = \ite(x < 0,\ 1-x,\ x+5), \qquad t = a + y,
\]
and then observes that the branch \texttt{t > \#0} depends on this conditional
value. The formula is small, but it already demonstrates the mechanism:
merging reduces the number of states while introducing a conditional value that
enters the next branch condition.

The final value \texttt{r} at \texttt{Lend} gives the contrasting case. An
\(\ite\)-representation can be constructed for that value as well, but there is
no later branch condition after that point. This candidate is less dangerous
with respect to future branch queries: the conditional value exists, but it is
not pushed into another condition.

The important distinction is future use. The local existence of an
\(\ite\)-expression is not enough to characterize merge cost; the relevant
question is whether the path-dependent value reaches a later condition.

\paragraph{Hash-consed expressions}
Besides path explosion, state merging also creates pressure on memory usage and
on the cost of traversing symbolic expressions. After a merge, the symbolic
store contains \(\ite\)-nodes. After SSA/phi-style representation and variable
versioning, the same subexpressions may appear in several merge candidates,
future-branch substitutions, and \(\widehat{Q}_{\mathrm{add}}\)-dependencies. If every step constructs
a fresh expression tree, identical fragments are duplicated even when they
represent the same object.

Here, hash-consing is background from prior symbolic-execution work: it shares
expression structure by assigning each expression node a structural key
consisting of the node kind, its local payload, and the identifiers of its
children \cite{veritesting}:

\[
\operatorname{key}(e) =
(\operatorname{kind}(e),\ \operatorname{payload}(e),\ \operatorname{id}(e_1),\ldots,\operatorname{id}(e_n)).
\]

The payload may contain an operator name, a constant, a variable name, a
variable version, or the information that the node is an \(\ite\). Before a new
node is created, the table of already constructed expressions is queried by
this key. If such a node already exists, its identifier is reused; otherwise, a
new node is allocated. As a result, symbolic expressions are stored as a shared
DAG rather than as a collection of independent trees.

For merged symbolic expressions, hash-consing reduces duplication and repeated
traversal without changing the meaning of the formulas. Counts of
\(\ite\)-nodes, expression size, phi-related merge candidates, versioned
dependencies, and future-branch substitution can all be computed over one
shared representation instead of over duplicated trees.

\section{Query Count Estimation and Dynamic State Merging}
\label{sec:qce}

\subsection{Query Count Estimation}
\label{sec:qce-theory}

Efficient State Merging introduces Query Count Estimation as a way to estimate
whether the expected benefit of merging outweighs the possible increase in the
complexity of future solver queries. In the source paper, \(Q_t\) estimates
the expected total number of future solver queries, \(Q_{\mathrm{add}}\)
concerns additional queries caused when a value that differs between states
becomes symbolic after merging, and \(Q_{\mathrm{ite}}\) concerns future
queries affected by conditional expressions introduced by the merge
\cite{efficient-merging}. Variable-specific estimates then support the
identification of hot variables.

The \(\ite\)-flow example in \Cref{sec:state-merging} illustrates the
\(Q_{\mathrm{ite}}\) side of this cost: a conditional value introduced at a
join reaches a later branch condition. Appendix~\ref{app:qce-manual-baseline}
manually applies the source-paper \(q\) recurrence to the MIR graph and
produces the oracle \(Q_t\) and \(Q_{\mathrm{add}}\) values used in the
comparison below.

The MIR experiment uses a separate hatted notation for the quantities computed
automatically from the CFG and recovered dependencies. For each block \(\ell\),
\(\widehat{Q}_t(\ell)\) is a CFG-derived estimate of future
conditional-branch pressure. For a variable \(v\),
\(\widehat{Q}_{\mathrm{add}}(\ell,v)\) is the part of that pressure attributed
to \(v\) by the selected dependency approximation. The local dependency
indicator \(C(\ell,v)\) is non-zero when the selected dependency approximation
decides that the branch condition at \(\ell\) depends on \(v\). The
experimental hot set is therefore

\[
\widehat{\operatorname{Hot}}(\ell)
=
\{v \mid
\widehat{Q}_{\mathrm{add}}(\ell,v)
> \alpha \cdot \widehat{Q}_t(\ell)\}.
\]

Conceptually, \(\beta\) controls how strongly a later branch influences
earlier locations. A larger value propagates more of the future pressure
backward through the CFG. The threshold \(\alpha\) controls how selective the
experimental hot-variable set is: small values mark more variables as hot,
while larger values keep only variables whose
\(\widehat{Q}_{\mathrm{add}}\) is a substantial fraction of
\(\widehat{Q}_t\).

In the default configuration, \(\alpha=0.35\) and \(\beta=0.50\). For an acyclic branch
graph, the MIR recurrence can be summarized as

\[
\widehat{Q}_t(\ell)=
\begin{cases}
1 + \beta \sum_{s \in succ(\ell)} \widehat{Q}_t(s), & \ell \text{ is conditional},\\
\sum_{s \in succ(\ell)} \widehat{Q}_t(s), & \text{otherwise},
\end{cases}
\]

and, for conditional blocks,

\[
\widehat{Q}_{\mathrm{add}}(\ell,v)=
C(\ell,v) + \beta \sum_{s \in succ(\ell)} \widehat{Q}_{\mathrm{add}}(s,v).
\]

For non-conditional blocks, the local \(C(\ell,v)\) term is omitted and only
successor pressure is propagated.

The same fixed-point computation gives both hatted quantities; only the local
cost function changes.

\begin{lstlisting}[style=pseudo,caption={Fixed-point propagation used for the hatted MIR estimates.},label={lst:qce-solver}]
algorithm SolveFutureCost(CFG, branch_successors, local_cost):
    q[block] = 0 for every block

    repeat until q changes by less than tolerance:
        next_q = empty vector

        for block in reverse block order:
            if block is a conditional branch:
                next_q[block] =
                    local_cost(block)
                    + beta * sum(q[s] for s in branch successors)
            else if block has one successor:
                next_q[block] = local_cost(block) + q[successor]
            else:
                next_q[block] =
                    local_cost(block)
                    + beta * sum(q[s] for s in successors)

        q = next_q

    return q
\end{lstlisting}

In the MIR example below, future conditional branches are used as the
query-pressure signal, while variable dependencies determine which source
values contribute to that pressure. For this acyclic graph, the automatic
\(\widehat{Q}_t\) recurrence has the same control-flow shape as the
source-paper \(Q_t\) oracle; the principal
experimental approximation is the recovered dependence relation used for
\(\widehat{Q}_{\mathrm{add}}\).

\subsection{Dependency Approximation Experiment}
\label{sec:qce-dependencies}

The MIR graph in \Cref{lst:mir-ite-preview} contains the dependency chain that
drives the QCE example. The later branch is written as \texttt{t > \#0};
\texttt{t} is computed as \texttt{a + y}; and \texttt{a} is defined on both
incoming paths from the earlier branch. Therefore, the source-level dependency
of the later branch is \(x,y\), even though the branch expression contains only
\texttt{t}.

The Appendix~\ref{app:qce-manual-baseline} oracle uses the source-paper
two-location relation \((\ell,v)\mathcal C(\ell',e)\): the value of \(v\) at
location \(\ell\) may affect the expression \(e\) evaluated at the later
location \(\ell'\). The four automatic dependency modes approximate this
relation from MIR syntax and CFG structure; the manual oracle is not a fifth
implemented mode.

The modes expose progressively deeper representations of the same dependency
chain rather than competing in an identical variable namespace. Syntactic mode
retains the branch temporary \texttt{t}; local expansion reaches \texttt{a}
and \texttt{y}; CFG propagation resolves \texttt{a} to the source variables
\texttt{x} and \texttt{y}; and versioned propagation records the same sources
with definition versions.

The syntactic approximation reads variables directly from branch expressions.
For \texttt{t > \#0}, it reports only \texttt{t}. This is a useful minimal
baseline because it locates query pressure cheaply: the analysis can already
see that block \(D\) contains a future branch. What it cannot explain is why
that branch is relevant to a merge performed earlier. The temporary
\texttt{t} is still a surface variable, not the source of the path-dependent
value.

The local approximation expands assignments inside the current block. Since
\texttt{t = a + y} appears immediately before the branch, local expansion
reports \texttt{a, y}. This removes the temporary \texttt{t} and exposes that
\texttt{a} and \texttt{y} are the values consumed by the later condition. The
important remaining gap is conceptual: the analysis knows that \texttt{a}
matters, but not why \texttt{a} differs between the incoming paths.

The CFG approximation carries dependency sets along CFG edges and unions them
at joins. It resolves \texttt{a} through both predecessor definitions and
therefore recovers \(x,y\), which matches the oracle source variables for this
example. This is the point where the earlier branch on \(x\) becomes connected
to the later branch on \(t\). The union at joins is path-insensitive, but for
this graph it is precise enough to explain the merge candidate.

The versioned approximation associates dependencies with source versions,
giving \(x_0,y_0\). Its numerical result is the same here because \(x\) and
\(y\) have one relevant source definition. The representation matters in
examples with redefinitions: it distinguishes a source value that reaches a
future branch from a later assignment to a variable with the same name. The
comparison is summarized in \Cref{tab:qce-dependency-modes}.

\subsection{QCE Results}
\label{sec:qce-results}

For the default configuration \((\alpha=0.35,\ \beta=0.50)\),
Table~\ref{tab:qce-comparison} places the automatic MIR estimates beside the
manually computed source-paper oracle. Table~\ref{tab:qce-comparison}(a)
aggregates the fixed-point result for each dependency mode. \textit{Vars} counts variables
with non-zero contribution, \textit{Blocks} counts blocks where future
dependency information is visible, \textit{Total \(\widehat{Q}_{\mathrm{add}}\)}
sums the variable-attributed pressure, and \textit{Hot refs} counts
block-variable references that pass the \(\widehat{\operatorname{Hot}}\)
threshold. The displayed hatted quantities, aggregate masses, and hot-reference
counts come from the automatic MIR analysis. Table~\ref{tab:qce-comparison}(b)
contains the manual oracle. \textit{Oracle \(L_1\)} is
computed separately by projecting those estimates onto the source-variable
domain of the manual oracle in Appendix~\ref{app:qce-manual-baseline}:

\[
L_{1,\mathrm{oracle}}(m)
=
\sum_{\ell}
\sum_{v\in V_{\mathrm{oracle}}}
\left|
\widehat{Q}_{\mathrm{add}}^{\,m}(\ell,v)
-
Q_{\mathrm{add}}(\ell,v)
\right|.
\]

For this example, \(V_{\mathrm{oracle}}=\{x,y\}\). Intermediate variables
reported by shallower modes, such as \texttt{t} or \texttt{a}, are not treated
as false positives; their effect appears as source-level attribution that has
not yet been recovered. For the versioned mode, source-version names such as
\(x_0\) and \(y_0\) are normalized to their corresponding source variables
before computing the distance. This article-level metric is distinct from the
dumper's internal \textit{\(L_1\Delta\) vs CFG} summary, which compares
automatic modes against CFG mode inside their reported variable spaces.

The main result is not the exact total mass but the recovered variable
identity. Syntactic mode retains only surface variables. Local mode improves
data-flow recovery inside the block, but still attributes relevance to the
intermediate join value \texttt{a}. CFG mode matches the oracle for this
example: it recovers the source variables \(x,y\), the per-block
\(Q_{\mathrm{add}}\) values, and the oracle hot-variable set. Versioned mode
has the same numerical agreement after normalizing \(x_0,y_0\) to \(x,y\).
Syntactic and local modes preserve part of the dependency at intermediate
program variables and therefore recover only part of the source-level
attribution measured by the oracle. Increasing \(\beta\) gives later branches
more influence over earlier blocks, but it does not change the qualitative
comparison between dependency approximations.

\subsection{Merge Decision}
\label{sec:qce-merge-decision}

The hatted estimates are then used for a lightweight merge-risk question:

\begin{quote}
If a join merges values from multiple predecessors, do those merged values flow
into future hot branch conditions?
\end{quote}

For each join candidate, the experiment identifies the join, recovers values
arriving from its predecessors, determines their source dependencies, compares
those dependencies with the experimental future-hot set
\(\widehat{\operatorname{Hot}}(\ell)\), and interprets the overlap. The
comparison uses the CFG dependency mode as the default source of future-hot
variables. The two relevant rows are shown in
\Cref{tab:state-merge-decision}.

There are two different sets involved in this comparison. The future-hot set
describes variables that can affect later branch pressure after the join. The
merged-source set describes variables behind values that actually differ
between the incoming states. A merge is risky when these sets overlap. A
future-hot variable that is not part of the path-specific difference is still
important to the later query, but it is not the reason this particular merge
introduces path-dependent structure.

The labels in \Cref{tab:state-merge-decision} are experimental labels for this
overlap test. The relevant comparison is
\[
\operatorname{DepsOfDifferingValue}
\cap
\widehat{\operatorname{Hot}}(\ell).
\]
A non-empty partial overlap is reported as \texttt{mixed}; a candidate after
which \(\widehat{Q}_t=0\) is reported as \texttt{no-future-query}.

The central row is \texttt{D}/\texttt{Ljoin}/\texttt{a}: branch-specific
definitions of \(a\) are reconstructed as
\(a=\ite(x<0,\ 1-x,\ x+5)\), the block computes \(t=a+y\), and the following
branch checks \(t>0\).

The merge of \texttt{a} at \texttt{Ljoin} has consequences beyond local SSA. The
merged value can influence the later branch through \texttt{t}. Both incoming
definitions of \(a\) depend on \(x\), so the path-specific difference being
merged is traced back to \(x\). The later branch depends on \(x\) and \(y\);
\(y\) is future-hot because it participates in \(t=a+y\), but it is not the
path-specific difference introduced by the merge. The relevant overlap is
therefore \(x\), which produces the \texttt{mixed} decision. The final-value
merge at \texttt{G}/exit is different: \texttt{r} may have different incoming
expressions, but no later condition consumes it, so no future branch pressure
remains and the row is labeled \texttt{no-future-query}.

The result supports the central state-merging point: a merge is potentially
risky not just because it creates an \(\ite\), but because the differing
value remains relevant to future branch conditions. This type of relevance
information could complement runtime merge heuristics without requiring the
full query model to be reconstructed statically.

A one-off call to \texttt{z3.simplify} was used to check the displayed form of
the reconstructed expression; the \(\ite\) structure remained unchanged
\cite{z3}.

\subsection{State Similarity and Dynamic State Merging}
\label{sec:dsm}

Dynamic State Merging uses QCE to define a similarity relation between states.
Two states at compatible program locations need not have identical symbolic
stores. If they differ only in variables that are cold with respect to future
queries, merging them is less likely to increase later formula cost. If they
differ in hot variables, merging is less attractive because those differences
are likely to appear in future branch conditions or other solver queries.

Thus, QCE supplies the relevance dimensions for comparing symbolic stores. The
state similarity question is not simply whether two stores are equal, but
whether they differ in variables that matter after the merge point. The
\texttt{Ljoin} example illustrates this relation: the path-dependent value of
\texttt{a} is relevant because it flows into \texttt{t > \#0}; a value that
differs only after the last branch is less important for future query
complexity.

A compact way to express the similarity test is

\[
\begin{aligned}
s_1 \sim_\ell s_2
\iff {} & \forall v \in \operatorname{Hot}(\ell): \\
& \mu_1(v)=\mu_2(v)
\lor \mathit{Sym}_I(\mu_1(v)) \\
& \lor \mathit{Sym}_I(\mu_2(v)).
\end{aligned}
\]

Here, \(\mathit{Sym}_I(e)\) means that expression \(e\) is already symbolic
with respect to the symbolic input set \(I\). Equality on hot variables is the
clean case. A difference may also be acceptable when at least one side is
already symbolic, because merging then does not newly turn that variable from a
concrete value into a symbolic value. Cold variables do not disappear from the
program state, but they do not dominate the similarity decision.

The MIR dependency experiment supplies the relevance part of this relation:
the hatted future-hot set indicates which variables should be compared.

DSM then addresses the scheduling side of the same problem. In ordinary dynamic
symbolic execution, active states are stored in a worklist, and
\texttt{pickNext} selects which state should be continued next. This driving
heuristic may optimize for depth, coverage, random-path behavior, or another
exploration criterion \cite{efficient-merging}. States that would be profitable
to merge may not naturally reach the same join point under that order.

DSM preserves the underlying driving heuristic but temporarily gives priority
to forwarding candidates. Predecessor history records recent states within a
bounded distance. When an active state is similar to a recent predecessor of
another active state, the algorithm treats this as a potential merge
opportunity and places suitable states into a forwarding set. These candidates
are not chosen because the normal search heuristic would necessarily prefer
them globally; they are chosen because the history relation suggests that they
may soon reach compatible locations. Forwarding ends naturally when states
merge, when the opportunity no longer holds, or when the similarity condition
stops applying. After that, exploration returns to the underlying search
strategy.

The dependency example illustrates the relevance information used by this
similarity relation, while DSM describes how such information is used during
exploration.

\section{Static Symbolic Execution and Veritesting}
\label{sec:veritesting}

Dynamic symbolic execution explores a program path by path: each state
corresponds to one branch history. Static Symbolic Execution uses a different
analysis unit. Instead of maintaining a separate state for each path, SSE tries
to construct a formula for an entire program fragment. For a small acyclic
region, this can be written in a simplified form as

\begin{equation}
\label{eq:sse-summary}
\operatorname{Summary}(R) =
\bigvee_{p \in \operatorname{paths}(R)}
PC_p \land \mathit{Update}_p.
\end{equation}

Here, \(PC_p\) describes the conditions required to follow path \(p\), and
\(\mathit{Update}_p\) describes the corresponding symbolic-store update. This makes SSE
a natural candidate for compressing local branch diamonds: several paths inside
a region can be represented by one summary structure in the form of
\Cref{eq:sse-summary}.

Veritesting uses this idea as part of dynamic analysis. At a candidate point,
the system recovers a local CFG from the current execution state, reduces that
CFG, identifies transition points, and applies SSE to the reduced acyclic
graph \cite{veritesting}. Branch contexts and symbolic stores are propagated
inside the static region; at joins they are merged into a summary rather than
materialized as one state per path. At transition points, DSE states are
recreated and ordinary path exploration resumes.

Transition points are the boundary between the two modes. They appear at exits
of the static region, at unresolved jumps, at calls or system interactions
that require the dynamic engine, and at loop-related cases handled through
unrolling, back-edge redirection, or an Incomplete Loop transition. This
alternation is what makes Veritesting different from using DSE or SSE alone:
static summaries are used locally, while the dynamic engine remains responsible
for the rest of the execution.

Java Ranger gives method calls a more central role because Java programs often
contain many small methods, and dynamically dispatched calls depend on runtime
receiver types \cite{java-ranger}. Inheritance and object-oriented dispatch
make method boundaries frequent and semantically important, so treating every
method call as an unconditional boundary would fragment multi-path regions.
Java Ranger therefore introduces Method Regions and Dynamic Method Region
Inlining: after runtime type information identifies the applicable target, a
method region can be instantiated and inlined into the regional summary.
Exceptional or otherwise unsuitable behaviors may remain separate Single-Path
Cases.

The MIR examples below use a narrower direct acyclic-region preview. The
preview asks:

\begin{quote}
Does this branch-shaped part of the CFG look like a small acyclic region that
could be summarized?
\end{quote}

The preview uses conditional entries, common exits or immediate
postdominators, the node set between them, back-edge presence, and directly
acyclic path sketches. It isolates one structural ingredient of Veritesting:
recognizing a region whose shape admits a summary.

The direct acyclic-region preview is shown in
\Cref{lst:acyclic-region-preview}.

\begin{lstlisting}[style=pseudo,caption={Direct acyclic-region preview.},label={lst:acyclic-region-preview}]
algorithm DirectAcyclicRegionPreview(CFG, branches, ipdom):
    for entry in branches:
        exit = ipdom(entry)
        region = nodes reachable before exit

        if region contains a cycle:
            row(entry, exit, region, status = "not directly acyclic")
        else:
            paths = bounded acyclic paths(entry, exit)
            row(entry, exit, region, paths)
\end{lstlisting}

The first test MIR graph contains two acyclic diamond structures
(\Cref{lst:mir-acyclic-diamond}).

\begin{lstlisting}[style=mirfull,caption={Acyclic diamond MIR input.},label={lst:mir-acyclic-diamond}]
param x
param y
ifTrue x > #0 goto Lpos
neg <-- -, #0, x
a <-- +, neg, y
goto Ljoin
Lpos: a <-- +, x, y
Ljoin: ifFalse y == #0 goto Lscale
r <-- +, a, #1
goto Lend
Lscale: r <-- *, a, y
Lend: return r
\end{lstlisting}

The preview identifies two regions that can be represented as compact acyclic
path sketches, as shown in \Cref{tab:region-preview-outcomes}.

Both regions have two syntactic branches and one common exit. This is the case
where a Veritesting-style approach is natural: instead of keeping two separate
DSE continuations, the local region can be viewed through a summary form.

The second test reuses the loop example from earlier sections
(\Cref{lst:mir-count-matches}). In that listing, line 5 is the loop header,
line 6 is the inner branch, and lines 10--11 form the latch and back edge.

Here there are two different candidate regions. One starts at the loop header
and contains a back edge. The other is the inner branch diamond inside the loop
body. The direct preview distinguishes these cases in
\Cref{tab:region-preview-outcomes}.

This behavior is important: the loop-shaped candidate is not directly acyclic
in the original CFG because it contains the back edge. A full Veritesting
treatment would require CFG reduction, unrolling, or a transition to dynamic
handling. The inner branch inside the loop is still a directly acyclic diamond
region and can be summarized separately.

A manual check gives the same classification; the same table records the
expected structural reason for each outcome.

Thus, the MIR examples illustrate a region-shape preview connected to
Veritesting. Directly acyclic branch diamonds are classified as summary-shaped
candidates, while a selected fragment with a back edge is marked as not
directly acyclic. The inner branch of the loop is still a directly acyclic
region, which shows that the preview depends on the selected CFG fragment,
rather than on whether the surrounding function contains a loop.

\section{Pattern-Based and Quantified State Merging}
\label{sec:quantified}

The previous section considered region-level summarization for acyclic CFG
fragments. Loops introduce a different kind of structure: not just several
branches inside one diamond region, but a repeated family of similar paths. If
the loop body contains a symbolic branch, ordinary DSE can represent different
combinations of branch outcomes across iterations. Standard state merging can
merge individual join points, but with many repetitions the resulting nested
\(\ite\)-expressions may still grow.

Quantified merging targets groups of execution-tree states whose histories do
not differ arbitrarily, but instead follow a regular pattern. The typical
source-paper shape is a family of leaves whose histories can be written as

\[
h(\pi(\mathit{leaf})) = \omega_1\omega_2^k\omega_3.
\]

The prefix \(\omega_1\) reaches the repeating fragment, \(\omega_2\) is the
loop-body pattern repeated \(k\) times, and \(\omega_3\) exits the fragment.
Compatible leaf states can then be partitioned into regular merging groups.
For such a group, repeated path constraints are encoded with a bounded
quantified schema whose \(k\) ranges over repetition counts represented by the
selected compatible states rather than every arbitrary unseen count; in the MIR
preview below, \(k\) is used as an explanatory loop-iteration index. Store
values are synthesized as functions of the repetition count \(k\) where
possible. Ordinary \(\ite\)-based merging remains available for variables whose
terms cannot be synthesized, and for groups where the regular pattern is not
applicable \cite{quantified-merging}.

The MIR experiment uses a CFG-level indexed-family preview. It approaches the
same regularity phenomenon from the CFG side: instead of constructing an
execution-tree merging group, it searches the loop for the structural elements
that would generate an indexed family: loop guard, induction update, inner
branch, guarded accumulator update, and repeat edge.

The loop/accumulator MIR example reuses \Cref{lst:mir-count-matches}. The
pattern experiment uses the loop guard on line 5, the inner branch on line 6, the
two \texttt{matches} updates on lines 7 and 9, and the induction/back-edge
shape on lines 10--11.

The corresponding source-level loop shape is the body of
\Cref{lst:c-count-matches}, especially lines 2--12.

The repeated inner branch is \(\mathit{target} = i\).

If the loop starts at \(i_0\) and exits at \(n\), the preview can write the
following compact indexed sketch:

\begin{equation}
\label{eq:quantified-matches}
\mathit{matches}_{\mathrm{out}}
=
\mathit{matches}_{\mathrm{in}}
+
\sum_{k=i_0}^{n-1}
\ite(\mathit{target} = k,\ 1,\ 0).
\end{equation}

\Cref{eq:quantified-matches} shows the regularity exposed by the preview: the
same conditional update is repeated over the loop index. The summation is an
explanatory sketch of the indexed family for this MIR loop. The loop guard
defines the range of \(k\), the induction update defines how \(k\) advances,
the inner branch defines the repeated condition, and the accumulator update
supplies the per-iteration contribution.

The compact recognition sketch in \Cref{lst:cfg-indexed-family-preview} checks
for the CFG components needed by this indexed view.

\begin{lstlisting}[style=pseudo,caption={CFG-level indexed-family preview.},label={lst:cfg-indexed-family-preview}]
algorithm CfgIndexedFamilyPreview(CFG, branches, assignments):
    find a loop back edge latch -> header
    recover nodes between header and latch

    guard = branch condition at header
    induction = linear update of the loop variable
    accumulator = inner branch with two constant updates

    if guard, induction, and accumulator are found:
        rewrite the branch condition with iteration variable k
        emit indexed sketch:
            accumulator += ite(condition(k), true_delta, false_delta)
\end{lstlisting}

The sketch looks for one back edge, recovers the loop body, and then identifies
three components: a guard, a linear induction update, and a guarded
accumulator. If all components are found, the result is an indexed sketch for
one iteration and for the range of \(k\).

For this MIR graph, the bounded node tree is shown in
\Cref{tab:quantified-node-tree}.

From this tree, the pattern-inference and manual-check summary in
\Cref{tab:quantified-pattern-inference} is extracted:

The same result can be described as the small automaton in
\Cref{lst:quantified-automaton}.

\begin{lstlisting}[style=pseudo,caption={Finite automaton view of the recognized loop pattern.},label={lst:quantified-automaton}]
S0 -- not(i < n) --------------------------> ACCEPT-EXIT
S0 -- i < n -------------------------------> S1
S1 -- target = i / matches += 1 -----------> S2
S1 -- not(target = i) / matches += 0 ------> S2
S2 -- i += 1; k += 1 ----------------------> S0
S0 + S1 + S2 recognized -------------------> ACCEPT
\end{lstlisting}

The automaton and the bounded node tree serve the same purpose: they show how a
parameterized family of repeated histories is generated from a small structural
pattern. The finite cases are loop exit, loop entry, true/false arm of the
inner branch, latch update, and repeat. Once these cases are recognized, the
repeated branch family can be described by one indexed schema.

A manual comparison with the target shape is included in the same table.

On the small loop example, the recovered guard, induction update, inner branch,
and accumulator update lead to an indexed sketch of the form

\[
\text{for}\; k \in [i_0,n) \colon
\quad
\mathit{matches} \mathrel{+\!=} \ite(\mathit{target} = k,\ 1,\ 0).
\]

Recognizing this structure before formula construction may help select cases
where quantified merging's execution-tree pattern search is likely to find a
regular group.

Automatic partial loop summarization is related, but it operates dynamically:
it observes loop executions, matches induction-variable and guard behavior,
predicts and confirms a symbolic iteration count, and emits partial pre-loop
and post-loop conditions. Quantified merging instead compresses regular groups
of symbolic states and path histories.


\section{Related Work}
\label{sec:related-work}

KLEE is the main baseline for the execution model used in this paper: it
executes LLVM IR over symbolic states, forks states at symbolic branches,
checks path feasibility with an SMT solver, and explores the resulting state
space~\cite{klee}. SymCC is related as an alternative implementation strategy
for symbolic execution, where instrumentation is compiled rather than
interpreted~\cite{symcc}.

Efficient State Merging is the closest reference for the QCE and DSM
discussion~\cite{efficient-merging}. In that work, \(Q_t\), \(Q_{\mathrm{add}}\), and
Dynamic State Merging guide real merge decisions during symbolic execution.
QCE supplies the future-relevance model used to reason about hot variables,
while DSM explains how state similarity can guide exploration toward merge
points. MultiSE is related through value summaries, which also represent
multiple path values inside one symbolic representation~\cite{multise}.

Veritesting motivates the direct acyclic-region preview~\cite{veritesting}. It
combines dynamic symbolic execution with static symbolic execution over
reduced local regions. Java Ranger extends a similar region-summary idea to
Java with Method Regions, Multi-Path Regions, dynamic method-region inlining,
Early-Returns Summarization, and exceptional Single-Path Cases
\cite{java-ranger}. The experiment in \Cref{sec:veritesting} examines one
region-shape question behind those approaches.

Quantified state merging motivates the CFG-level indexed-family
preview~\cite{quantified-merging}. The experiment in \Cref{sec:quantified}
examines whether a regular loop shape with a guard, induction update, inner
branch, and guarded accumulator update can expose an indexed family.

Loop summarization and compositional dynamic test generation provide broader
context for summarizing repeated or reusable behavior
\cite{loop-summarization,smart-summaries}. Heap-oriented symbolic execution
and modular heap summaries address richer memory structures than the scalar MIR
examples used in this paper~\cite{sep-logic-exec,heap-summaries}.
CFM-SE is a concrete targeted control-flow melding transformation for reducing
branching pressure~\cite{cfm-se}. Targeted program transformations for
symbolic execution argue more broadly that such transformations should be
first-class scalability components~\cite{targeted-transforms}.

\section{Conclusion}
\label{sec:conclusion}

State merging trades state count for formula complexity. The examples in this
paper show why future relevance matters: a conditional value is most important
when it flows into later branch conditions, and dependency precision determines
which variables appear hot. In the MIR example, CFG-level propagation is
enough to recover the relevant source variables, while purely
syntactic or block-local approximations lose part of the explanation.

Hot variables connect QCE to Dynamic State Merging by defining the dimensions
used for state similarity. Veritesting generalizes the merging idea from values
at a join point to acyclic CFG regions, and quantified merging generalizes it
again to regular loop-path families. Some of the required information resembles
data already computed by control-flow and dependency analyses, but the main
lesson is about state merging itself: merging is useful when the path
differences that remain in formulas are controlled and future-relevant
differences are handled carefully.


\bibliographystyle{IEEEtran}
\bibliography{seminar-paper}

\clearpage
\onecolumn

\appendices
\renewcommand{\theHtable}{appendix.\Alph{section}.\arabic{table}}
\section{Tables}
\label{app:tables}
\begingroup
\renewcommand{\tighttable}{\footnotesize\renewcommand{\arraystretch}{0.95}}
\setlength{\tabcolsep}{3pt}
\setlength{\floatsep}{4pt plus 1pt minus 1pt}
\setlength{\textfloatsep}{4pt plus 1pt minus 1pt}
\setlength{\intextsep}{4pt plus 1pt minus 1pt}
\setlength{\abovecaptionskip}{2pt}
\setlength{\belowcaptionskip}{2pt}

\begin{table}[H]
\caption{Dependency sets recovered by the four MIR dependency approximations.}
\label{tab:qce-dependency-modes}
\centering
\tighttable
\begin{tabularx}{\textwidth}{@{}lllY@{}}
\toprule
Mode & Direct branch variable & Recovered dependency & Note \\
\midrule
Syntactic
& \texttt{t}
& \texttt{t}
& Reads variables directly present in the condition. \\
Local
& \texttt{t}
& \texttt{a}, \texttt{y}
& Expands the assignment \texttt{t = a + y}. \\
CFG
& \texttt{t}
& \texttt{x}, \texttt{y}
& Propagates dependencies through predecessor blocks. \\
Versioned
& \texttt{t}
& \texttt{x0}, \texttt{y0}
& Records source definitions explicitly. \\
\bottomrule
\end{tabularx}
\end{table}

\begin{table}[H]
\caption{Automatic MIR estimates and manual source-paper QCE oracle.}
\label{tab:qce-comparison}
\centering
\begin{minipage}[t]{0.58\textwidth}
\centering
{\footnotesize (a) Automatic MIR estimates}
\vspace{1pt}

\centering
\tighttable
\setlength{\tabcolsep}{2.5pt}
\begin{tabular}{@{}lrrrrrr@{}}
\toprule
Mode & Vars & Blocks & \(\max\widehat{Q}_t\) & \(\Sigma\widehat{Q}_{\mathrm{add}}\) & Oracle \(L_1\) & Hot \\
\midrule
Syntactic & 2 & 2 & 2.00 & 7.00 & 10.00 & 7 \\
Local & 3 & 2 & 2.00 & 12.00 & 5.00 & 12 \\
CFG & 2 & 2 & 2.00 & 12.00 & 0.00 & 10 \\
Versioned & 2 & 2 & 2.00 & 12.00 & 0.00 & 10 \\
\bottomrule
\end{tabular}
\end{minipage}
\hfill
\begin{minipage}[t]{0.38\textwidth}
\centering
{\footnotesize (b) Manual source-paper QCE oracle}
\vspace{1pt}

\centering
\tighttable
\setlength{\tabcolsep}{2.5pt}
\begin{tabular}{@{}lrrrr@{}}
\toprule
Blocks & \(\beta\) & \(Q_t\) & \(Q_{\mathrm{add}}(x)\) & \(Q_{\mathrm{add}}(y)\) \\
\midrule
Entry, \(A\) & 0.50 & 2.00 & 2.00 & 1.00 \\
Entry, \(A\) & 0.75 & 2.50 & 2.50 & 1.50 \\
Entry, \(A\) & 0.90 & 2.80 & 2.80 & 1.80 \\
\(B\), \(C\), \(D\) & any & 1.00 & 1.00 & 1.00 \\
\(E\), \(F\), \(G\), Exit & any & 0.00 & 0.00 & 0.00 \\
\bottomrule
\end{tabular}
\end{minipage}
\end{table}

\begin{table}[H]
\caption{Merge-candidate overlap with experimental future-hot variables.}
\label{tab:state-merge-decision}
\centering
\tighttable
\begin{tabularx}{\textwidth}{@{}llllllY@{}}
\toprule
Join & Var & Source deps & Future hot & Overlap & Decision & Reason \\
\midrule
\texttt{D}/\texttt{Ljoin}
& \texttt{a} & \texttt{x} & \texttt{x}, \texttt{y} & \texttt{x}
& mixed
& \texttt{a} flows into \texttt{t = a + y}; the later branch uses \texttt{t}. \\
\texttt{G}/exit
& \texttt{r} & \texttt{x}, \texttt{y} & none & none
& no-future-query
& No later branch pressure remains after this join. \\
\bottomrule
\end{tabularx}
\end{table}

\begin{table}[H]
\caption{Direct acyclic-region preview results.}
\label{tab:region-preview-outcomes}
\centering
\tighttable
\begin{tabularx}{\textwidth}{@{}lllY Y@{}}
\toprule
Case & Entry--exit & Nodes & Observed outcome & Structural check \\
\midrule
First diamond
& \texttt{A}\(\to\)\texttt{D} & \texttt{A,B,C}
& 2 paths: \(\lnot(x > 0) \rightarrow D\), \(x > 0 \rightarrow D\)
& Acyclic branch with common exit. \\
Second diamond
& \texttt{D}\(\to\)\texttt{G} & \texttt{D,E,F}
& 2 paths: \(y = 0 \rightarrow G\), \(\lnot(y = 0) \rightarrow G\)
& Acyclic branch with common exit. \\
Loop header
& \texttt{B}\(\to\)\texttt{G} & \texttt{B,C,D,E,F}
& Not directly acyclic; no path sketch.
& \texttt{F}\(\to\)\texttt{B} is a back edge. \\
Inner branch
& \texttt{C}\(\to\)\texttt{F} & \texttt{C,D,E}
& 2 paths: \(\mathit{target} = i \rightarrow F\), \(\lnot(\mathit{target} = i) \rightarrow F\)
& Loop-internal, but the selected region is acyclic. \\
\bottomrule
\end{tabularx}
\end{table}

\begin{table}[H]
\caption{Bounded node tree for the CFG-level indexed-family preview.}
\label{tab:quantified-node-tree}
\centering
\tighttable
\begin{tabularx}{\textwidth}{@{}rllllY@{}}
\toprule
Depth & Node & Via & Next & Role & Detail \\
\midrule
0 & \texttt{B} & loop header & \texttt{C} / \texttt{G}
& loop guard & branch on \(i < n\) \\
1 & \texttt{C} & \(i < n\) & \texttt{D} / \texttt{E}
& inner branch & branch on \(\mathit{target} = i\) \\
2 & \texttt{D} / \texttt{E} & true / false arm & \texttt{F}
& accumulator update & \texttt{matches += 1} or \texttt{matches += 0} \\
3 & \texttt{F} & latch & \texttt{B}
& induction update & \texttt{i += 1}; back edge to the loop header \\
\bottomrule
\end{tabularx}
\end{table}

\begin{table}[H]
\caption{CFG-level indexed-family components recovered from the MIR example.}
\label{tab:quantified-pattern-inference}
\centering
\tighttable
\begin{tabularx}{\textwidth}{@{}lY Y Y@{}}
\toprule
Pattern & Evidence & Inferred sketch & Manual check \\
\midrule
Loop guard
& \texttt{B}
& \(i < n\)
& Matches the loop range guard. \\
Induction update
& \texttt{F}
& \texttt{i = + i \#1}; \(i := i + 1\)
& Linear constant-step update. \\
Guarded accumulator
& \texttt{D} / \texttt{E}
& \texttt{matches += ite(target = k, 1, 0)}
& Matches the per-iteration branch/update shape. \\
Indexed sketch
& loop \(\texttt{B} \leftarrow \texttt{F}\)
& \texttt{for k in [i0,n): matches += ite(target = k,1,0)}
& Range shape matches the indexed summation in the text. \\
\bottomrule
\end{tabularx}
\end{table}
\endgroup

\clearpage

\section{Manual QCE Oracle from Efficient State Merging}
\label{app:qce-manual-baseline}

This appendix manually applies the original QCE recurrence from Efficient
State Merging to \Cref{lst:mir-ite-preview} \cite{efficient-merging}. It uses
a manually determined source-level dependence relation and produces the
unhatted \(Q_t\) and \(Q_{\mathrm{add}}\) oracle used to evaluate the automatic
MIR approximations in \Cref{sec:qce-results}.

The source recurrence is

\[
q(\ell,c)=
\begin{cases}
\beta q(\operatorname{succ}(\ell),c)
+\beta q(\ell',c)
+c(\ell,e),
& \operatorname{instr}(\ell)=\operatorname{if}(e)\operatorname{goto}\ell',
\\
0, & \operatorname{instr}(\ell)=\operatorname{halt},
\\
q(\operatorname{succ}(\ell),c), & \text{otherwise}.
\end{cases}
\]

The two quantities used here are

\[
Q_t(\ell)=q\left(\ell,\lambda(\ell',e).1\right)
\]

and

\[
Q_{\mathrm{add}}(\ell_0,v)
=
q\left(
\ell_0,
\lambda(\ell',e).
\ite\bigl((\ell_0,v)\mathcal C(\ell',e),1,0\bigr)
\right).
\]

The relation \((\ell_0,v)\mathcal C(\ell',e)\) states that the value of
variable \(v\) at location \(\ell_0\) may affect expression \(e\) evaluated at
the later location \(\ell'\).

The CFG successor relation is:
Entry\(\to A\); \(A\) falls through to \(B\) and branches to \(C\);
\(B\to D\); \(C\to D\); \(D\) falls through to \(E\) and branches to \(F\);
\(E\to G\); \(F\to G\); \(G\to\)Exit; Exit is halt. The graph is acyclic, so
no loop bound is needed.

There are two conditional expressions. For \(e_A=(x<0)\), the value of \(x\)
at Entry and \(A\) may affect \(e_A\). For \(e_D=(t>0)\), the source-level
semantics before the branch is

\[
t = \ite(x < 0,\ 1 - x,\ x + 5) + y.
\]

Therefore, the value of \(x\) and the value of \(y\) at any of Entry, \(A\),
\(B\), \(C\), and \(D\) may affect \(e_D\). No later conditional expression is
reachable from \(E\), \(F\), \(G\), or Exit.

At the final conditional block,

\[
Q_t(D)=1,\qquad
Q_{\mathrm{add}}(D,x)=1,\qquad
Q_{\mathrm{add}}(D,y)=1.
\]

Blocks \(B\) and \(C\) deterministically reach \(D\), so the same values
propagate without a \(\beta\) factor:

\[
Q_t(B)=Q_t(C)=1,
\quad
Q_{\mathrm{add}}(B,x)=Q_{\mathrm{add}}(C,x)=1,
\quad
Q_{\mathrm{add}}(B,y)=Q_{\mathrm{add}}(C,y)=1.
\]

At block \(A\), the first branch contributes one local query, and both
successors contribute with factor \(\beta\):

\[
Q_t(A)=1+\beta(Q_t(B)+Q_t(C))=1+2\beta.
\]

For \(x\), the branch at \(A\) contributes locally and the later branch at
\(D\) contributes through both successors:

\[
Q_{\mathrm{add}}(A,x)
=
1+\beta(Q_{\mathrm{add}}(B,x)+Q_{\mathrm{add}}(C,x))
=
1+2\beta.
\]

For \(y\), there is no local contribution at \(A\), but the later condition at
\(D\) depends on \(y\):

\[
Q_{\mathrm{add}}(A,y)
=
\beta(Q_{\mathrm{add}}(B,y)+Q_{\mathrm{add}}(C,y))
=
2\beta.
\]

Entry deterministically reaches \(A\), so Entry has the same values as \(A\).
For \(E\), \(F\), \(G\), and Exit, no future conditional expression remains:
\(Q_t=0\) and \(Q_{\mathrm{add}}(\ell,x)=Q_{\mathrm{add}}(\ell,y)=0\). The
resulting oracle values for \(\beta \in \{0.50,\ 0.75,\ 0.90\}\) are reported
in Table~\ref{tab:qce-comparison}(b), alongside the automatic MIR estimates in
Table~\ref{tab:qce-comparison}(a).

\end{document}
